%% =====================================================================
%%  dictamen.tex  --  Plantilla de DICTAMEN de comision (MDD)
%% ---------------------------------------------------------------------
%%  La estructura la fija el esquema `dictamen` del ecosistema (Reglamento
%%  del Congreso art. 70; Manual de Tecnica Legislativa 2021 sec. XI) y el
%%  compilador la exige. Variantes de formato: anexos A-F del Manual
%%  (aprobacion en sus terminos, texto sustitutorio, no aprobacion...).
%% =====================================================================
\documentclass{dictamen}
%% =====================================================================
%%  metadata.tex  --  DATOS del dictamen
%% =====================================================================
\setmeta{numero}{001}
\setmeta{anio}{2026}
% \setmeta{fecha}{11 de setiembre del 2026}   % comentado: usa la fecha de hoy
\setmeta{firmante}{\meta{autor}}   % por defecto firma el titular del perfil
\setmeta{cargofirmante}{Diputado de la Rep\'ublica}
\setmeta{comision}{Comisi\'on de [Inserta aqu\'i]}
\setmeta{periodo}{Per\'iodo Anual de Sesiones 2026-2027}

\esquemacondicion{aprobacion}    % dictamen de aprobacion: con formula legal (quitar si no aprueba)

\begin{document}

\componente{sumilla}
\sumilla{[Inserta aqu\'i: en una frase, qu\'e propone el proyecto y qu\'e recomienda la comisi\'on.]}

\componente{encabezado}
\begin{center}
  {\fuentedisplaysemi\color{colortitulos}\MakeUppercase{Dictamen recaído en el Proyecto de Ley N.° [Inserta aquí]}}\\[2pt]
  \meta{comision} \textperiodcentered\ \meta{periodo}
\end{center}
\vspace{0.6\baselineskip}

\componente{presentacion}
\noindent Ha ingresado para dictamen de la \meta{comision} el Proyecto de Ley N.°
[Inserta aqu\'i], presentado por [Inserta aqu\'i], que propone [Inserta aqu\'i].

\componente{situacion}
\componente{procedimentales}
[Fecha de presentaci\'on, decreto de env\'io, comisiones dictaminadoras.]
\componente{parlamentarios}
[Proyectos anteriores sobre la misma materia, por n\'umero de expediente y estado.]

\componente{contenido}
[Qu\'e propone el proyecto, art\'iculo por art\'iculo.]

\componente{marco}
\componente{nacional}
[Constituci\'on, leyes y reglamentos aplicables, verificados sobre el texto.]

\componente{opiniones}
[Sumilla de las opiniones ciudadanas y t\'ecnicas recibidas (art. 70.2 y 70.3).]

\componente{analisis}
\componente{tecnico}
[Inserta aqu\'i.]
\componente{necesidad}
[Inserta aqu\'i.]
\componente{efecto}
[Normas cuya derogaci\'on o modificaci\'on se propone (art. 70.5).]
\componente{analisis-opiniones}
[Inserta aqu\'i.]
\componente{costo-beneficio}
[Efectos monetarios y no monetarios, impacto econ\'omico y presupuestal, sectores (art. 70.6).]

\componente{conclusion}
Por las consideraciones expuestas, la \meta{comision}, de conformidad con el
literal b) del art\'iculo 70 del Reglamento del Congreso de la Rep\'ublica,
recomienda la \textbf{APROBACI\'ON} del Proyecto de Ley N.° [Inserta aqu\'i],
con el siguiente texto sustitutorio:

\componente{formula}
\rotulolegislativo{Texto sustitutorio}
\titulodelaley{Ley que [Inserta aqu\'i]}
\articulo{1}{Objeto de la ley}
[Inserta aqu\'i.]

\componente{estilo}
\formuladeestilo

\componente{firmas}
\firmaslegislativas

\end{document}
